\section{Introduction}
\label{sec:introduction}

A growing fraction of the public web is generated by large language
models, and the next generation of foundation models will train on a
substrate that is partly their own ancestors' output. When a single
generative model is trained on its own outputs in a closed loop, the
distribution it represents narrows: tails disappear first, modes
collapse, and the model converges to a low-variance regime
\citep{shumailov2024nature, alemohammad2023mad, dohmatob2024tale}. This
phenomenon, \emph{model collapse}, has been replicated across modalities,
architectures, and training regimes.

\para{Why this matters.}
Almost every published collapse study fixes the unit of analysis at a
\emph{single} self-training model. The web, however, is populated by
\emph{many} LLMs from different vendors, with different pretraining
priors, system prompts, and retrieval pipelines. The policy-relevant
question is not ``does one model collapse on its own outputs'' --- the
answer to that is yes --- but rather: how many distinct LLM agents,
sharing a common substrate, are enough to keep that substrate
\emph{healthy}? Borrowing the term from population ecology, is there a
\emph{minimum viable population} (\mvp)
\citep{shaffer1981minimum} of LLMs below which collapse is unavoidable?
And what kind of differentiation between agents counts as ``different
individuals''?

\para{Gap.}
Three recent ecosystem-level studies each leave a piece of the question
open. \citet{hodel2026epistemic} test only $M \in \{1, 2, 4, 16\}$
copies of one base model differentiated by training-data segment, find
that the \emph{optimal} $M$ grows monotonically with iteration count, but
never report an \mvp at any fixed horizon. \citet{wang2025web} run a
network of $N{=}3$ pretrained-different LLMs sharing a retrieval pool
but never sweep $N$. \citet{vu2025recursive} treat $N{=}2$ model families
theoretically but do not test $N{>}2$. None of these works compare
multiple axes of inter-agent differentiation, and none compare the
fine-tuning ecosystem to the cheaper retrieval-augmented ecosystem on
identical metrics.

\para{Our approach.}
We close all three gaps in one experimental program. We run two
complementary harnesses on the \emph{same} diversity metric suite
(perplexity, distinct bigrams, mean pairwise embedding distance, and
the Vendi-score Hill--Shannon Diversity, hereafter \hsd). E1 is a
fine-tuning ecosystem where $N$ copies of \distilgpt retrain on each
others' \wikitext outputs in replace mode. E2 is a retrieval-augmented
ecosystem where $N$ frozen API LLMs each query a shared, growing
``internet'' of their own posts. E3 fixes $N{=}4$ in the retrieval
harness and varies four axes of individuality: random seed, retrieval
data slice, system-prompt persona, and pretrained model family.

\para{Quantitative preview.}
We find that \emph{$N{=}2$ frozen API LLMs from different pretraining
families} preserve full diversity for $T{=}12$ iterations
(slope of mean pairwise distance: $-0.002 \pm 0.004$, $p = 0.44$),
while in the fine-tuning harness even $N{=}16$ agents collapse by a
factor of $2.10\times$ in perplexity over only $T{=}6$ iterations.
Across the four diversity axes at $N{=}4$, model-family diversity
preserves $\mathbf{1.9\times}$ the terminal inter-agent distance of the
next-best axis (persona or data slice). The largest single $N$-doubling
benefit in the fine-tuning ecosystem is \mbox{$N{=}8 \to N{=}16$}, which
cuts terminal perplexity by $26\%$.

\para{Contributions.}
\begin{itemize}[leftmargin=*,itemsep=0pt,topsep=0pt]
    \item We introduce a population-ecology framing of recursive
    LLM training, formalising the \mvp question and giving the first
    explicit \mvp estimates for both fine-tuning and retrieval-based
    ecosystems on a common metric suite.
    \item We extend the only published $N$-vs-collapse sweep
    \citep{hodel2026epistemic} from $N \le 16$ to a controlled
    statistical comparison across $N \in \{1, 2, 4, 8, 16\}$ in the
    fine-tuning regime and $N \in \{1, 2, 3, 5, 8\}$ in the retrieval
    regime, holding compute roughly constant.
    \item We run the first head-to-head comparison of four orthogonal
    axes of inter-agent differentiation (seed, data slice, prompt
    persona, model family) at fixed $N{=}4$, and show that
    \emph{architecture diversity is the only axis that buys
    collapse-resistance}.
    \item We reconcile the divergent answers from the two harnesses
    into a single governance-relevant claim: it is the
    \emph{training-time} link, not the inference-time loop, that drives
    ecosystem collapse.
\end{itemize}

\para{Paper organization.}
\secref{sec:related} positions the work against the single-model and
ecosystem-level collapse literature. \secref{sec:methodology} details
the two harnesses, metrics, and statistical analysis.
\secref{sec:results} reports E1, E2, and E3 results, with \mvp
estimates and axis ranking. \secref{sec:discussion} discusses
limitations and what the cross-method gap means for AI deployment.
\secref{sec:conclusion} closes with bottom-line answers and
recommended follow-ups.
